\chapter{Diseño e Implementación}
\label{cap:implementacion}

Este capítulo constituye el núcleo técnico de mi Trabajo de Fin de Grado. Aquí describo cómo se ha construido el sistema, qué partes son desarrollo propio y qué partes reutilizan herramientas ya existentes, y cómo se han interconectado los distintos microservicios para convertir una colección de componentes sueltos en un asistente corporativo funcional. Para no recargar la lectura, el detalle de implementación más extenso ---fragmentos de código, contratos de clase y flujos paso a paso--- se traslada a los Anexos~\ref{anexo:pipeline}--\ref{anexo:finetuning}, y aquí se mantiene la visión de arquitectura y las decisiones de diseño.

% ============================================================
\section{Filosofía de Diseño: Necesidad, Solución, Iteración}
\label{sec:filosofia}
% ============================================================

El sistema no se diseñó como un bloque monolítico, sino que emergió de un proceso iterativo de identificación de necesidades y acoplamiento progresivo de soluciones. Cada nueva funcionalidad surgió como respuesta a una limitación concreta detectada durante el uso del sistema:

\begin{enumerate}
    \item Necesidad inicial: Un chat con un LLM local. \textit{Solución:} Ollama + OpenWebUI \cite{ollama2024docs,openwebui2024docs}. El sistema comenzó como una instalación básica de OpenWebUI conectada a Ollama, proporcionando una interfaz de chat con modelos locales.

    \item Primera limitación: El modelo respondía de memoria, sin acceso a documentos internos. \textit{Solución:} Se desarrolló un backend FastAPI con pipeline RAG y base de datos vectorial Qdrant. Esto exigió un mecanismo para interceptar los mensajes del usuario antes de que llegaran directamente al modelo: el Pipeline de OpenWebUI.

    \item Segunda limitación: Los documentos cambiaban frecuentemente en SharePoint y la ingesta manual era inviable. \textit{Solución:} Se creó el servicio Indexer con sincronización automática vía Microsoft Graph API \cite{msgraph2024docs}.

    \item Tercera limitación: Muchos PDFs corporativos eran escaneados y no contenían texto seleccionable. \textit{Solución:} Se integró PaddleOCR con aceleración GPU en el Indexer.

    \item Cuarta limitación: Los usuarios necesitaban acceder a información que no estaba en los documentos internos, como noticias, páginas web o normativa. \textit{Solución:} Se añadieron los modos de búsqueda web, scraping y consulta del BOE. En la interfaz principal, la consulta normativa se resuelve mediante API; además, se desarrolló un servidor MCP como línea de integración futura.

    \item Quinta limitación: Sin visibilidad sobre el rendimiento del sistema, era imposible detectar cuellos de botella. \textit{Solución:} Se implementó el stack de observabilidad: métricas Prometheus en el backend, NVIDIA DCGM para GPU, dashboards Grafana y logs centralizados con Loki/Promtail \cite{prometheus2026docs,dcgmExporter2026docs,grafana2026docs}.

    \item Sexta limitación: Las respuestas idénticas saturaban la GPU innecesariamente. \textit{Solución:} Se interpuso LiteLLM como proxy con caché en Redis.

    \item Séptima limitación: El modelo no usaba la terminología corporativa adecuada. \textit{Solución:} Se exploró el fine-tuning con LoRA sobre Qwen 2.5.

    \item Octava limitación: Los analistas de datos necesitaban consultar tablas de datos estructurados sin aprender SQL. \textit{Solución:} Se desarrolló un agente de traducción NL$\rightarrow$SQL (\texttt{core/sql\_agent.py}) que genera y ejecuta consultas \texttt{SELECT} seguras sobre PostgreSQL mediante LiteLLM, con validación por lista blanca y auto-corrección ante errores de sintaxis.

    \item Novena limitación: El control de acceso mediante cabeceras HTTP era insuficiente para entornos multi-tenant exigentes. \textit{Solución:} Se implementó validación criptográfica de tokens JWT de Azure AD (\texttt{core/auth.py}), con descarga dinámica de claves públicas JWKS, activable con \texttt{AZURE\_JWT\_VALIDATION=true}.

    \item Décima limitación: La expansión de consulta en el endpoint de búsqueda lanzaba el cálculo en paralelo pero nunca recogía el resultado: estaba efectivamente deshabilitada en producción sin que los tests lo detectaran. \textit{Solución} (v2.2.0): se añadió la recogida asíncrona del futuro, se ejecutan búsquedas con las consultas expandidas y se incorporan sus resultados sin duplicados; además se limpiaron imports perezosos y logs de depuración en \texttt{core/retrieval.py}.

    \item Undécima limitación: El pipeline RAG usaba siempre el mismo modelo y los mismos parámetros (\texttt{top\_k=5}) para cualquier consulta. \textit{Solución} (v2.3.0): se desarrolló el \texttt{QueryProcessor} con \textit{intent detection} (clasifica en FACTUAL/PROCEDURAL/ANALYTICAL/CONVERSATIONAL y ajusta \texttt{top\_k}), \textit{smart model routing} (analíticas a \texttt{qwen2.5-32b}; el resto a JARVIS) y \textit{query expansion} (reformula y fusiona resultados antes del reranker). El \texttt{MemoryManager} pasó a usar \texttt{qwen2.5-32b} para summarizar conversaciones largas.
\end{enumerate}

Esta aproximación iterativa ---identificar una carencia, evaluar alternativas, integrar la solución y validar el resultado--- se repitió a lo largo de los siete sprints del proyecto. Cada nueva pieza se acopló al ecosistema respetando la separación de responsabilidades y las interfaces HTTP estándar entre servicios, permitiendo sustituir cualquier componente sin afectar al resto.

% ============================================================
\section{Arquitectura General del Sistema}
\label{sec:arquitectura-general}
% ============================================================

El sistema se compone de 15 servicios principales y varios servicios auxiliares de soporte y arranque, todos contenerizados con Docker, orquestados mediante Docker Compose \cite{dockerCompose2026docs} y aislados en una red virtual propia (\texttt{rag-network}). El flujo resumido de tecnologías se presentó en el Capítulo~\ref{cap:tecnologias}; aquí se detalla la función de cada servicio dentro de la arquitectura funcional.

\begin{table}[H]
\centering
\caption{Servicios principales de la arquitectura y su función.}
\label{tab:servicios-completos}
\begin{tabular}{llc}
\toprule
\textbf{Servicio} & \textbf{Función} & \textbf{Puerto} \\
\midrule
NGINX              & Proxy inverso, terminación TLS        & 8443 (externo) \\
OpenWebUI          & Interfaz de usuario conversacional     & 8080 (interno) \\
Pipelines (JARVIS) & Agente encaminador inteligente         & 9099 (interno) \\
Backend (FastAPI)  & Motor RAG, scraping, búsqueda, OCR     & 8000 (interno) \\
Indexer (FastAPI)  & Sincronización SharePoint, ingesta     & 8001 (interno) \\
Ollama             & Ejecución local de LLMs                & 11434 (interno) \\
LiteLLM            & Proxy unificado de modelos             & 4000 (interno) \\
Qdrant             & Base de datos vectorial                & 6333 (interno) \\
PostgreSQL         & Base de datos relacional               & 5432 (interno) \\
Redis              & Caché de respuestas                    & 6379 (interno) \\
MCP-BOE Server     & Consulta del BOE                       & 8010 (interno) \\
MinIO              & Almacenamiento de objetos (backup)     & 9000 (interno) \\
Prometheus         & Recolección de métricas del sistema    & 9090 (interno) \\
Grafana            & Dashboards de observabilidad           & 3000 (interno) \\
pgAdmin            & Administración de PostgreSQL           & 5050 (interno) \\
\bottomrule
\end{tabular}
\end{table}

Todos los servicios comparten una red Docker interna de tipo \textit{bridge} denominada \texttt{rag-network}. Únicamente NGINX expone un puerto al exterior (8443), garantizando que ningún servicio interno es accesible directamente desde fuera del servidor. Los servicios que requieren aceleración GPU (Ollama, PaddleOCR) la reciben mediante la directiva \texttt{deploy.resources.reservations} de Compose y el \texttt{nvidia-container-toolkit} del host; los servicios más intensivos declaran además límites de CPU y memoria (\texttt{deploy.resources.limits}) para que ninguno monopolice los recursos del host.

En la validación final (26 de mayo de 2026), el despliegue convivía con otros servicios Docker en el mismo host. Para evitar colisiones de puertos se mantuvieron los puertos internos y se publicaron accesos locales diferenciados (p.~ej. OpenWebUI en \texttt{localhost:3002}, Backend en \texttt{localhost:8002}, Qdrant en \texttt{localhost:6335}). Esta separación no modifica la arquitectura interna, pero permite ejecutar el sistema completo sin detener otros servicios del entorno de desarrollo.

\begin{figure}[H]
\centering
\includegraphics[width=0.95\textwidth]{cap_tfg_docker_stack.png}
\caption{Estado del stack Docker en la validación final. La captura muestra los servicios principales de \texttt{TFG-RAG-Clean}, su estado de ejecución y los puertos alternativos publicados para la demo local.}
\label{fig:docker-stack-final}
\end{figure}

\begin{figure}[H]
\centering
\includegraphics[width=0.9\textwidth]{cap_tfg_backend_swagger.png}
\caption{Documentación interactiva del Backend FastAPI en el entorno local de validación: API activa, endpoints de scraping, búsqueda y salud, y contrato OpenAPI expuesto.}
\label{fig:backend-swagger-final}
\end{figure}

\subsection{Flujo de Datos End-to-End}

Para ilustrar cómo se conectan las piezas, se describe el recorrido completo de una consulta RAG típica:

\begin{enumerate}
    \item El \textbf{usuario} escribe su pregunta en el navegador (OpenWebUI).
    \item \textbf{NGINX} recibe la petición HTTPS y la reenvía internamente a OpenWebUI (puerto 8080).
    \item \textbf{OpenWebUI} envía el mensaje al pipeline JARVIS mediante HTTP POST al puerto 9099.
    \item \textbf{JARVIS} analiza el texto, detecta la intención RAG, obtiene los departamentos autorizados del usuario y construye una petición al Backend.
    \item El \textbf{Backend} genera el embedding de la consulta con SentenceTransformer, ejecuta una búsqueda híbrida en Qdrant y recupera los $k$ fragmentos más relevantes.
    \item El Backend compone un \textit{prompt} estructurado con el contexto recuperado y lo envía a LiteLLM.
    \item \textbf{LiteLLM} verifica si existe una respuesta cacheada en Redis. Si no, delega al modelo en Ollama.
    \item \textbf{Ollama} ejecuta la inferencia en GPU y devuelve los tokens generados en streaming.
    \item La respuesta fluye de vuelta: LiteLLM $\rightarrow$ Backend $\rightarrow$ JARVIS $\rightarrow$ OpenWebUI $\rightarrow$ NGINX $\rightarrow$ navegador.
\end{enumerate}

Todo el flujo se completa típicamente en menos de 3.5 segundos, con el primer token visible para el usuario en menos de 1.2 segundos gracias al streaming.

\begin{figure}[H]
\centering
\includegraphics[width=0.85\textwidth]{screenshot_openwebui_jarvis.png}
\caption{Interfaz de OpenWebUI con el modelo JARVIS seleccionado. El usuario escribe su consulta y recibe la respuesta en streaming.}
\label{fig:screenshot-openwebui}
\end{figure}

% ============================================================
\section{Desarrollo Propio vs. Herramientas Externas}
\label{sec:propio-vs-externo}
% ============================================================

Una pregunta fundamental en cualquier proyecto de integración es: ¿qué parte del trabajo es desarrollo original y qué parte es configuración de herramientas existentes? La Tabla~\ref{tab:reuso-desarrollo} clasifica cada componente del sistema según su origen.

\begin{table}[H]
\centering
\caption{Resumen sintético de reutilización y desarrollo dentro del TFG.}
\label{tab:reuso-desarrollo}
\resizebox{\textwidth}{!}{%
\begin{tabular}{p{3.0cm}ccccp{5.2cm}}
\toprule
\textbf{Componente} & \textbf{Reutilizado} & \textbf{Extendido} & \textbf{Desarrollado en el TFG} & \textbf{Integrado} & \textbf{Papel en el sistema} \\
\midrule
OpenWebUI                  & \checkmark & \checkmark & $\times$   & \checkmark & Interfaz conversacional web, persistencia de chats y base para autenticación corporativa. \\
Pipeline JARVIS            & $\times$   & $\times$   & \checkmark & \checkmark & Encaminamiento de intenciones, memoria de sesión y conexión entre interfaz y backend. \\
Backend RAG                & $\times$   & $\times$   & \checkmark & \checkmark & Recuperación documental, scraping, búsqueda web, API BOE y lógica principal del sistema. \\
Indexer                    & $\times$   & $\times$   & \checkmark & \checkmark & Sincronización con SharePoint, OCR e ingesta automatizada de documentos. \\
Ollama y LiteLLM           & \checkmark & \checkmark & $\times$   & \checkmark & Inferencia local, alias de modelos, caché y desacoplo entre servicios y modelos. \\
Qdrant                     & \checkmark & \checkmark & $\times$   & \checkmark & Almacenamiento vectorial y recuperación documental con filtros por colección. \\
Servidor MCP-BOE           & $\times$   & $\times$   & \checkmark & \checkmark & Consulta normativa estructurada y prueba de integración MCP. \\
Prometheus, Grafana y DCGM & \checkmark & \checkmark & $\times$   & \checkmark & Observabilidad operativa del backend, la base de datos y la GPU. \\
\bottomrule
\end{tabular}%
}
\end{table}

En resumen, el trabajo de desarrollo propio abarca más de 6.000 líneas de código Python distribuidas entre el Pipeline JARVIS (1.498 líneas), el Backend ($\sim$3.000), el Indexer ($\sim$1.200) y el servidor MCP-BOE ($\sim$400), además de toda la configuración declarativa del ecosistema Docker (Compose, LiteLLM, NGINX, Prometheus y los dashboards de Grafana). Las herramientas externas proporcionan las capacidades de infraestructura (base de datos, motor LLM, OCR, navegador), pero la lógica de negocio, la orquestación y la inteligencia del sistema son enteramente originales.

% ============================================================
\section[Pipeline de OpenWebUI]{El Mecanismo Pipeline de OpenWebUI (Hook)}
\label{sec:hook-pipeline}
% ============================================================

El Pipeline de OpenWebUI es el punto de conexión fundamental entre la interfaz de usuario y todo el backend del sistema \cite{openwebui2024docs}. Sin este mecanismo, OpenWebUI sería simplemente una interfaz para chatear directamente con Ollama. El Pipeline es lo que convierte esa interfaz genérica en un asistente corporativo inteligente.

En ingeniería de software, un \textit{hook} es un punto de extensión que permite inyectar código personalizado en un flujo de ejecución predefinido sin modificar el sistema original. En OpenWebUI, el \textit{hook} se materializa en el sistema de Pipelines: cuando un usuario envía un mensaje, en lugar de reenviarlo directamente a Ollama, OpenWebUI lo envía primero al servicio Pipeline, que procesa el mensaje, ejecuta la lógica necesaria y devuelve la respuesta. Desde la perspectiva de OpenWebUI, el Pipeline \textit{es} el modelo.

La conexión se establece por configuración: el Pipeline se despliega como contenedor que expone una API HTTP en el puerto 9099, y OpenWebUI se configura con \texttt{OPENAI\_API\_BASE\_URLS} apuntando a él. Al arrancar, OpenWebUI descubre el Pipeline por su nombre (\texttt{JARVIS}) y lo muestra como un modelo seleccionable. El detalle de implementación ---el contrato de la clase \texttt{Pipeline}, el método \texttt{pipe()} con respuesta en \textit{streaming} (\texttt{yield}) y el sistema \textit{Valves} de configuración dinámica desde la interfaz de administración--- se recoge en el Anexo~\ref{anexo:pipeline}.

% ============================================================
\section{Agente Encaminador JARVIS: Inteligencia Central}
\label{sec:jarvis-detalle}
% ============================================================

El agente JARVIS es la pieza de software más compleja del proyecto, con 1.498 líneas de código Python. Su función es analizar cada mensaje del usuario y decidir a cuál de los modos de operación derivarlo. La detección de intenciones se implementa en el método \texttt{\_detect\_intent()}, que evalúa el mensaje contra un conjunto de reglas heurísticas (coincidencia de palabras clave y expresiones regulares) en un \textit{orden de prioridad estricto}: primero las condiciones más específicas (archivos adjuntos, imágenes, URLs) y después las más genéricas (keywords de RAG, búsqueda web, conversación). La Tabla~\ref{tab:intenciones-detalle} resume las categorías, su prioridad y los criterios de activación.

\begin{table}[H]
\centering
\caption{Categorías de intención y su prioridad de evaluación.}
\label{tab:intenciones-detalle}
\begin{tabular}{clp{7cm}}
\toprule
\textbf{Prioridad} & \textbf{Modo} & \textbf{Criterio de activación} \\
\midrule
0 & Decisión pendiente   & Existe una decisión de URL esperando respuesta del usuario \\
1 & Archivo adjunto      & Se detecta un PDF o documento en el cuerpo de la petición \\
2 & Visión/OCR           & Imagen adjunta en el mensaje actual \\
3 & Referencia a imagen  & Mención a ``la imagen'' con imagen previa en memoria \\
4 & Referencia a web     & Mención a ``la web'' o ``la página'' con contenido web en memoria \\
5 & Ayuda                & Keywords: ``qué puedes hacer'', ``ayuda'', ``help'' \\
6 & Listar documentos    & Commands: ``/docs'', ``/listar''; keywords: ``qué documentos tienes'' \\
7 & BOE / Legal          & Keywords: ``busca en el BOE'', ``ley de'', ``real decreto'' \\
8 & RAG Documental       & Keywords: ``documento'', ``política'', ``manual'', ``procedimiento'' \\
9 & URL detectada        & Expresión regular que detecta \texttt{https://...} en el mensaje \\
10 & Búsqueda web        & Keywords: ``busca en internet'', ``noticias'', ``precio'' \\
11 & Estado del sistema  & Keywords: ``cómo va'', ``estado'', ``indexación'' \\
12 & Sticky Mode         & Seguimiento contextual del modo de la pregunta anterior \\
13 & Chat general        & Default: cualquier mensaje que no active otro modo \\
\bottomrule
\end{tabular}
\end{table}

\begin{figure}[H]
\centering
\includegraphics[width=0.85\textwidth]{cap_chat_lista_docs.png}
\caption{Respuesta real del modo de listado documental. JARVIS consulta el estado del indexador y devuelve los documentos disponibles, junto con la información reciente de ingesta.}
\label{fig:documentos-disponibles}
\end{figure}

Más allá del encaminamiento, JARVIS incorpora tres mecanismos de contexto: el \textit{Sticky Mode} (mantiene el modo de la pregunta anterior ante consultas de seguimiento breves y ambiguas), la \textbf{memoria de sesión por usuario} (recuerda la última imagen y la última web analizadas, y gestiona decisiones pendientes) y el \textbf{control de acceso departamental} (mapea los grupos de Azure AD del usuario a colecciones de Qdrant y filtra la búsqueda). El detalle de estos tres mecanismos, con ejemplos de código, se recoge en el Anexo~\ref{anexo:pipeline}.

% ============================================================
\section{Backend RAG: Motor de Búsqueda Semántica}
\label{sec:backend-rag}
% ============================================================

El Backend constituye el motor central de procesamiento del sistema. Implementado en FastAPI, fue objeto de una refactorización significativa: inicialmente era un único fichero \texttt{main.py} de más de 1.900 líneas que se dividió en una arquitectura modular con paquetes separados para esquemas (\texttt{schemas/}), servicios (\texttt{services/}), endpoints (\texttt{api/}), procesamiento (\texttt{processing/}) e integraciones (\texttt{integrations/}).

La \textbf{ingesta} de un documento sigue seis etapas: recepción (API, Indexer o scraper), detección de tipo (texto nativo vs. escaneado), extracción de texto (\texttt{pdfplumber}, PaddleOCR o la cadena Trafilatura/Readability/BeautifulSoup), \textit{chunking} semántico (fragmentos de 512--1024 tokens con 10--15\% de solapamiento), generación de \textit{embeddings} (vector denso de 384 dimensiones con \path{paraphrase-multilingual-MiniLM-L12-v2}) y almacenamiento en Qdrant con metadatos. El procedimiento completo se detalla en el Anexo~\ref{anexo:backend}.

\begin{figure}[H]
\centering
\includegraphics[width=0.9\textwidth]{cap_tfg_qdrant_dashboard.png}
\caption{Estado de las colecciones en Qdrant durante la demo final. La colección \texttt{documents\_TFGDEMO} ---documentación sintética y segura del propio TFG--- aparece separada del resto, con 13 puntos y configuración híbrida dense+sparse.}
\label{fig:qdrant-tfgdemo}
\end{figure}

\subsection{Búsqueda Híbrida}
\label{subsec:busqueda-hibrida}

La recuperación de contexto combina tres etapas complementarias:

\begin{enumerate}
    \item \textbf{Búsqueda densa (HNSW)}: la consulta se transforma en un embedding de 384 dimensiones y se buscan los fragmentos más próximos en el espacio vectorial mediante el índice HNSW de Qdrant.
    \item \textbf{Búsqueda dispersa (BM25)}: se complementa con vectores dispersos generados con \texttt{Qdrant/bm25} (vía \texttt{fastembed}), que favorece coincidencias léxicas exactas ---identificadores, códigos, siglas o artículos legales--- que la búsqueda semántica podría no capturar.
    \item \textbf{Reordenación con reranker}: los candidatos se fusionan y se reordenan con un \textit{Cross-Encoder} (\texttt{cross-encoder/ms-marco-MiniLM-L-6-v2}), que evalúa conjuntamente el par (consulta, fragmento) para obtener una relevancia más precisa antes de seleccionar los $k$ fragmentos finales (típicamente $k=5$).
\end{enumerate}

El \textit{prompt} enviado al LLM combina un \textit{system prompt} restrictivo ---``responde únicamente con el contexto proporcionado, no inventes''--- con una temperatura baja ($0.1$--$0.2$), lo que elimina prácticamente las alucinaciones en modo RAG. La construcción exacta del \textit{prompt} se muestra en el Anexo~\ref{anexo:backend}.

% ============================================================
\section{Extracción Web, OCR e integraciones}
\label{sec:extraccion-integraciones}
% ============================================================

El sistema añade varias fuentes de información más allá de los documentos internos, cuyo detalle de implementación se recoge en los anexos:

\begin{itemize}
    \item \textbf{Scraping y búsqueda web} (Anexo~\ref{anexo:backend}): un motor con Playwright renderiza páginas dinámicas (JavaScript) y una cadena de \textit{fallback} (Trafilatura $\rightarrow$ Readability $\rightarrow$ BeautifulSoup) extrae el contenido principal. La búsqueda en internet usa \texttt{duckduckgo\_search} con \textit{fallback} a scraping HTML.
    \item \textbf{OCR} (Anexo~\ref{anexo:backend}): PaddleOCR con aceleración GPU ---y \textit{fallback} automático a CPU--- procesa los PDFs escaneados, con paralelización mediante Ray durante la ingesta.
    \item \textbf{SharePoint} (Anexo~\ref{anexo:integraciones}): el Indexer sincroniza cada 5 minutos mediante MSAL y Microsoft Graph, con estrategia \textit{delta} incremental para no reprocesar documentos.
\end{itemize}

% ============================================================
\section{Integración con el Boletín Oficial del Estado (BOE)}
\label{sec:boe-impl}
% ============================================================

La integración con el BOE se implementa mediante un servidor MCP (\textit{Model Context Protocol}) construido con \texttt{fastmcp}, que transforma peticiones en lenguaje natural en llamadas estructuradas a la API de Datos Abiertos del BOE. Esta integración no es una búsqueda web genérica, y merece una justificación específica: en el tipo de empresa para la que se diseña el sistema, el acceso a normativa legal (legislación laboral, protección de datos, contratación pública) es frecuente y tiene consecuencias directas. Un error en la interpretación de una norma puede traducirse en incumplimientos o sanciones.

El BOE presenta características que lo distinguen como caso de uso propio: las leyes se identifican por códigos oficiales (p.~ej. \texttt{BOE-A-2018-16673}) y no por títulos que el usuario recuerde; un mismo texto se referencia coloquialmente de muchas formas (``la LOPD'', ``la ley de privacidad''); la información relevante suele estar en artículos numerados concretos; y tanto la exactitud como la fuente son cruciales. Por ello se implementó una capa de resolución semántica que traduce el lenguaje del usuario al identificador normativo correcto antes de invocar la API oficial. El mecanismo de mapeo y la discusión sobre la integración MCP vs.\ API directa se detallan en el Anexo~\ref{anexo:integraciones}.

% ============================================================
\section{Seguridad Perimetral y Autenticación}
\label{sec:seguridad-impl}
% ============================================================

La seguridad del sistema se organiza en capas. \textbf{NGINX} actúa como primera barrera y único servicio expuesto al exterior: termina el TLS (los microservicios internos hablan en texto plano dentro de la red aislada), hace \textit{proxy pass} preservando cabeceras de trazabilidad y habilita el \textit{upgrade} a WebSocket necesario para el streaming. La autenticación de usuarios se delega en \textbf{Azure AD vía OpenID Connect (OIDC)}: los usuarios entran con sus credenciales corporativas y los atributos del token (nombre, correo, grupos) alimentan el control de acceso departamental.

\begin{figure}[H]
\centering
\includegraphics[width=0.85\textwidth]{cap_openwebui_login_actual.png}
\caption{Pantalla real de autenticación de OpenWebUI en el despliegue local: inicio de sesión clásico y acceso mediante SSO corporativo.}
\label{fig:login-sso}
\end{figure}

Además de la autenticación perimetral, el Backend implementa una segunda capa de \textbf{validación criptográfica de tokens JWT} (\texttt{core/auth.py}) para clientes externos y herramientas MCP: descarga las claves públicas JWKS del \textit{tenant}, verifica cada token Bearer y mapea los grupos a colecciones autorizadas, devolviendo \texttt{HTTP 403} si el acceso no es válido. Se activa con \texttt{AZURE\_JWT\_VALIDATION=true}. Finalmente, el \textbf{aislamiento de red} (\texttt{rag-network}, tipo \textit{bridge}) impide el acceso externo directo a PostgreSQL, Redis, Qdrant u Ollama. El flujo JWT paso a paso se detalla en el Anexo~\ref{anexo:integraciones}.

% ============================================================
\section{Observabilidad, caché y almacenamiento}
\label{sec:observabilidad-impl}
% ============================================================

La monitorización respondió a una necesidad real: durante las primeras fases no existía visibilidad sobre el rendimiento, la saturación de la GPU ni los patrones de uso, y la optimización era intuitiva en lugar de basada en datos. El stack adoptado combina \textbf{Prometheus} (métricas del backend: latencias y percentiles, throughput, distribución de modos, errores), \textbf{NVIDIA DCGM Exporter} (VRAM, utilización, temperatura, potencia de la GPU), \textbf{Grafana} (tres dashboards: backend, GPU y base de datos) y \textbf{Loki + Promtail} (agregación de logs de todos los contenedores, consultables junto a las métricas en Grafana) \cite{prometheus2026docs,grafana2026docs,dcgmExporter2026docs}. El detalle de métricas, telemetría y dashboards está en el Anexo~\ref{anexo:integraciones}.

\begin{figure}[H]
\centering
\includegraphics[width=0.9\textwidth]{cap_tfg_grafana_jarvis_dashboard.png}
\caption{Dashboard principal de Grafana durante la validación final: estado operativo del backend, Qdrant y Prometheus, además de métricas de memoria, CPU y ficheros abiertos.}
\label{fig:grafana-final}
\end{figure}

Esta observabilidad condujo a decisiones concretas de optimización: se validó que $\sim$35\% de las consultas repetidas se beneficiaban de la caché, se ajustó el valor de $k$ al detectar que más de 5 fragmentos no mejoraban la respuesta, y se reprogramó la ingesta OCR a horarios de baja actividad al ver que monopolizaba la GPU.

En cuanto al rendimiento y el almacenamiento, \textbf{Redis} actúa como caché integrada con LiteLLM: almacena un hash de cada consulta junto con su respuesta (TTL de 1 hora), con una tasa de aciertos observada de $\sim$35\% que reduce la carga sobre la GPU. \textbf{MinIO} proporciona almacenamiento de objetos compatible con S3 que preserva una réplica aislada de los documentos binarios ingeridos, permitiendo reindexar desde cero sin volver a las fuentes originales (SharePoint, web).

% ============================================================
\section{Agente SQL: Consultas en Lenguaje Natural sobre Datos Estructurados}
\label{sec:sql-agent}
% ============================================================

El pipeline RAG estándar opera solo sobre texto no estructurado, pero muchos sistemas empresariales guardan información crítica en tablas relacionales (métricas, inventarios, datos financieros). Para consultarlas sin saber SQL se desarrolló el agente SQL (\texttt{backend/app/core/sql\_agent.py}), que sigue tres pasos: (1)~\textbf{inspección del esquema} de PostgreSQL para inyectarlo como contexto; (2)~\textbf{generación} de una consulta \texttt{SELECT} con LiteLLM, instruyendo al modelo a producir solo lecturas; y (3)~\textbf{ejecución con auto-corrección}, reintentando con el mensaje de error como contexto (hasta dos veces). La consulta se valida contra una lista blanca (solo \texttt{SELECT}/\texttt{WITH}/\texttt{EXPLAIN}) y se ejecuta con \textit{timeout}; los mecanismos de seguridad se detallan en el Anexo~\ref{anexo:backend}. El endpoint unificado \texttt{POST /api/v1/query} enruta automáticamente entre RAG y SQL según la naturaleza de la consulta, de forma transparente para el usuario.

% ============================================================
\section{Fine-Tuning con LoRA}
\label{sec:finetuning-impl}
% ============================================================

Para adaptar el modelo al dominio corporativo se exploró el ajuste fino con LoRA (\textit{Low-Rank Adaptation}), que inserta matrices de rango reducido en las capas de atención sin modificar los pesos originales, reduciendo los parámetros entrenables a una fracción y haciendo viable el entrenamiento con una sola GPU. El modelo resultante (\texttt{rag-qwen-ft:latest}) es el que sirve de base a JARVIS. El procedimiento completo ---generación del dataset, entrenamiento y evaluación--- se detalla en el Anexo~\ref{anexo:finetuning}.

Una conclusión relevante del proyecto es que, para tareas de RAG puras, el \textit{fine-tuning} del LLM resultó \textit{menos} decisivo que un buen \textit{prompt} restrictivo combinado con una recuperación de calidad (chunking, embeddings, búsqueda híbrida y reranker). El ajuste fino sí aportó valor en la coherencia terminológica y en la adherencia al estilo de respuesta corporativo, pero la inteligencia del sistema descansa sobre todo en la calidad del contexto recuperado, no en la modificación de los pesos del modelo.
